% Requires \usepackage{booktabs} and \usepackage{tabularx}. Spans both columns (table*).
\begin{table*}[ht]
  \centering
  \caption{Structural decomposition of the six detectors (backbone / neck / head / post-processing).}
  \label{tab:detector_anatomy}
  \footnotesize
  \setlength{\tabcolsep}{4pt}
  \renewcommand{\arraystretch}{1.3}
  \begin{tabularx}{\textwidth}{l *{6}{>{\centering\arraybackslash}X}}
    \toprule
     & YOLOv9-m & YOLOv10-m & YOLO11-m & YOLO12-m & YOLO26-m & RT-DETR-L \\
    \midrule
    Backbone
     & GELAN (CNN)
     & CSPNet + PSA
     & C3k2 + C2PSA
     & R-ELAN (area-attn.)
     & C3k2-based CNN
     & HGNetv2 \\
    \addlinespace
    Neck
     & GELAN PAN-FPN
     & PAN-FPN
     & PAN-FPN w/ C3k2
     & PAN-FPN (A$^2$ blocks)
     & PAN-FPN
     & Hybrid encoder (AIFI + CCFM) \\
    \addlinespace
    Head
     & Decoupled AF (+PGI aux., train-only)
     & Decoupled dual (1-to-1 + 1-to-many)
     & Decoupled AF
     & Decoupled AF
     & End-to-end, DFL-free
     & Transformer decoder + IoU-aware query sel. \\
    \addlinespace
    Post-proc.
     & NMS
     & None
     & NMS
     & NMS
     & None
     & None \\
    \bottomrule
  \end{tabularx}
  \par\smallskip
  \footnotesize ``AF'' = anchor-free; ``PGI aux.'' = programmable-gradient auxiliary branch, active only during training; ``A$^2$'' = area-attention. The five YOLO variants share the single-stage backbone--neck--head anatomy and differ in their constituent blocks; RT-DETR-L replaces the neck with a hybrid encoder (intra-scale AIFI + cross-scale CCFM) and the head with a transformer decoder, giving NMS-free set prediction. YOLO26-m and RT-DETR-L details should be confirmed against the trained checkpoints.
\end{table*}
